\documentclass[10pt,a4paper]{article}


\usepackage[utf8]{inputenc}
\usepackage{amsmath}
\usepackage{graphicx}
\usepackage{geometry}
\usepackage{cite}
\usepackage{caption}
\geometry{margin=1cm}

\title{\textbf{BTEP Project Proposal}}
\date{\today}
\author{}

\begin{document}

\maketitle

\section*{Project Title}
Design and Hardware Implementation of a Real-Time Image Tracking Software using Verilog HDL

\section*{Student Information}
\begin{itemize}
    \item \textbf{Name:} Ayush Purohit
    \item \textbf{Roll Number:} AU2240179
    \item \textbf{Program:} B.Tech Computer Science and Engineering
    \item \textbf{Supervisor name:} Professor Mazad Zaveri
\end{itemize}

\section*{Introduction}
Image tracking is a basic issue in computer vision which entails the estimation of the movement and condition of an object in a 
sequence of successive frames of a video stream. It is important in the use of the system in video surveillance, autonomous vehicles, robotics, 
human-computer interaction, and intelligent transportation systems. Most conventional image tracking algorithms are generally run in software on a 
general-purpose processor (CPU) or graphics processing unit (GPU). Though flexible they tend to be unable to fulfill real-time requirements, especially in 
embedded or power-constrained systems. \\

\noindent Image processing tasks can be broken down to fine-grain parallelism (exploited by designers via Hardware Description Languages (HDLs) like Verilog) through 
hardware implementation. Filtering, state updates and matrix arithmetic are some of the operations that can be implemented simultaneously and make 
significant enhancements in throughput over the sequential execution of software. Moreover, deterministic timing behavior of hardware makes such 
implementation very suitable to real-time systems. \\

\noindent The reason I chose this topic is I am very much interested in hardware architecture, how algorithms can be converted into a series of hardware units like adders, 
multipliers, shifters etc., and have wanted to do something in the vector processing domain, specifically image processing. As such, image tracking particularly 
interested me as I believe it will be a good challenge for me as a project, while simultaneously having several real world applications in various sectors.

\section*{Problem Statement}
Image tracking algorithms, usually implemented through general purpose processor, or a graphics processing unit, tend to be slow or sometimes unable to fulfill real-time 
requirements. Can those algorithms be designed and translated to direct hardware implementations? If so, with what benefits or limitations?

\section*{Objectives of the Study}
\textbf{Main Objective:} To design and implement a real-time image tracking algorithm on hardware using Verilog HDL. \\
\textbf{Specific Objectives:} 
\begin{enumerate}
    \item To simulate said hardware on an FPGA board.    
    \item To measure the delay and number of clocks required to perform the operation

\end{enumerate}

\section*{Scope of the Project}
The project will cover the design and hardware implementation of the core computational components of an image tracking system using Verilog HDL. It will emphasize more on 
architecture design, and real-time feasability of image tracking algorithm in hardware. 

\noindent In scope will be:
\begin{enumerate}
    \item Selection of Tracking Algorithm that is Hardware Friendly
    \item Hardware Design using Verilog
    \item Fixed Point Arithmetic Implementation
    \item Simulation and Functional Verification
    \item Performance Evaluation
    \item Design for FPGA Implementation
\end{enumerate}

\\

\noindent Not in scope will be:
\begin{enumerate}
    \item Advanced Computer Vision Techniques
    \item Floating-Point Hardware Design
    \item End to End Vision System
    \item Software Integration
\end{enumerate}


\section*{Methodology}
\begin{enumerate}
    \item The project focuses on analytical and experimental methodology. Therefore, it will use synthetic test data, generate reference results baseline using 
    software models from MATLAB or Python, analyze standard image tracking algorithm to identify and extract mathematical models to identify hardware suitable operations 

    \item Tools used will be:
    \begin{itemize}
        \item Icarus Verilog and GTKWave for HDL Simulation
        \item FPGA Synthesis and analysis tools
        \item Software tools for reference modelling
        \item Other support tools like IDE, git etc.
    \end{itemize}

    \item Steps to be followed
        \begin{enumerate}
            \item Algorithm selection and Modelling
            \item Fixed-Point Representation Design
            \item Hardware Design
            \item Module Development
            \item Function Simulation and Verification
            \item Performance Evaluation
            \item Analysis and Documentation
        \end{enumerate}
\end{enumerate}

\section*{Expected Outcomes}
The project will make the feasibility and effectiveness of implementing an image tracking algorithm in hardware with Verilog, 
and obtain a real time performance with reasonable accuracy and resource use. Particularly:
\begin{enumerate}
    \item Understanding Image Tracking Algorithm in Hardware Context
    \item Practical Experience in Verilog HDL
    \item Fixed Point Arithmetic and Analysis Skills
    \item Performance and Resource Evaluation
    \item Architechtural Optimization Insights
    \item Feasability of Hardware Based Tracking
\end{enumerate}

\section*{Timeline}
Below is the \textbf{tentative} timeline of the project. \\

\noindent\textbf{Phase I:} Literature Review and Problem Definition \\
\textbf{Phase II:} Algorithm Modelling and Fixed Point Analysis \\
\textbf{Phase III:} Hardware Architecture Design \\
\textbf{Phase IV:} Verilog Implementation and Integration \\
\textbf{Phase V:} Verification and Performance Evaluation \\
\textbf{Phase VI:} Documentation and Final Presentation \\

\textbf{Tentative} Gantt Chart is shown in Figure 1

\begin{figure}[h]
    \centering
    \includegraphics[width=0.7\textwidth]{"./Gantt Tentative.jpeg"}
    \label{Gantt_Chart}
    \caption{\textbf{Tentative} Gantt Chart}

\end{figure}

\nocite{*}
\bibliographystyle{ieeetr}
\bibliography{./refs}

\end{document}
